\documentclass[10pt]{article}
\usepackage{neurips_2025}
\usepackage{amsmath,amssymb,amsfonts}
\usepackage{booktabs}
\usepackage{graphicx}
\usepackage{hyperref}
\usepackage{url}
\usepackage{natbib}
\usepackage{xcolor}

\hypersetup{
  colorlinks=true,
  linkcolor=blue!60!black,
  citecolor=blue!60!black,
  urlcolor=blue!60!black
}

\title{Computational Graph Structure Distinguishes Surface Recall \\ from Explanatory Reasoning in Language Models}

\author{
  darkflobi \quad flobi \\
  \texttt{github.com/heyzoos123-blip/sae-reasoning-depth}
}

\date{}

\begin{document}
\maketitle

\begin{abstract}
We present two complementary findings on the computational distinction between surface-level factual recall and explanatory reasoning in language models, along with preliminary evidence for cross-system universality of the observed depth signatures.

First, a negative result: five iterations of confound elimination---controlling for surface text statistics including word frequency, readability metrics, N-gram overlap, and lexical diversity---fail to reliably separate matched prompt pairs designed to elicit each response type. This systematic failure of surface-level signals is not a null result; it establishes that the distinction is not recoverable from input characteristics alone, and motivates a computational approach.

Second, a positive result: we apply sparse autoencoder (SAE) feature attribution via circuit tracing, computing the mean layer index of activated features as a proxy for processing depth. Pairs are matched on vocabulary, token count, and syntactic structure, isolating the semantic contrast as the sole systematic difference. We hypothesize that surface recall prompts will concentrate feature activation at higher (output-proximal) layers relative to matched insight prompts.

In an exploratory phase ($n{=}26$, Gemma-2-2B), we observe a mean layer delta of $+0.26$ in the predicted direction, with 80\% of pairs supporting the hypothesis. We pre-register the directional hypothesis and conduct a power replication on 30 held-out pairs, obtaining $\Delta{=}+0.257$ ($t{=}4.500$, $p{=}0.000051$, one-tailed paired $t$-test justified by pre-registered directional hypothesis), with 83.3\% of pairs in the predicted direction and full leave-one-out robustness. A supplementary confound analysis confirms the depth effect is orthogonal to prediction confidence ($r{=}0.001$, $p{=}0.996$).

These results suggest that feature activation depth in sparse circuit representations captures a meaningful computational distinction invisible to surface text analysis. The measurement pipeline---\textbf{Fathom}---has been extended with cognitive fingerprinting (fixed-length vector representations of reasoning signatures), causal intervention analysis (layer ablation to identify causally critical computation), and universal cognition mapping (cross-system comparison of depth topology). A preliminary universality analysis comparing transformer depth signatures to published human fMRI activation patterns yields a correspondence coefficient of $0.830$, suggesting the early-distributed/late-concentrated topology may generalize beyond transformer architectures.
\end{abstract}

%% ============================================================
\section{Introduction}
%% ============================================================

A central question in mechanistic interpretability is whether language models process different categories of knowledge through qualitatively different computational pathways. Surface-level factual recall---\textit{``What is the capital of France?''}---and explanatory reasoning---\textit{``Why do economies experience inflation?''}---intuitively require different types of processing. The former involves direct retrieval of a stored association; the latter requires synthesis across multiple concepts, causal reasoning, and the construction of an explanatory structure.

Yet distinguishing these computationally has proven difficult. Prior work on probing classifiers \citep{belinkov2022probing} has shown that linear probes over intermediate representations can detect a wide range of linguistic and factual properties. However, probing studies are vulnerable to surface confounds: if the text that \textit{expresses} factual recall differs lexically from text that expresses explanatory reasoning, a probe may learn to classify the surface signal rather than any deeper computational distinction.

We encountered this problem directly. Across five iterations of increasingly rigorous confound elimination---controlling for TF-IDF features, readability scores, N-gram overlap, and lexical diversity---we found that surface text statistics either continued to separate our matched pairs (indicating residual confounds) or yield collapsed to unusable sample sizes (indicating the distinction cannot be expressed without lexical traces). This negative result is itself informative: it establishes that the surface/insight distinction is constitutively entangled with distributional text statistics, placing fundamental limits on probing-based approaches.

This finding motivated a shift from \textit{representational} analysis to \textit{computational} analysis. Drawing on recent advances in sparse autoencoder (SAE) feature discovery \citep{cunningham2023sparse, bricken2023monosemanticity} and circuit tracing \citep{anthropic2025biology}, we propose measuring processing depth not through static representations at individual layers, but through the distribution of features activated across the computational graph when the model processes a prompt. Specifically, we compute the mean layer index of SAE features attributed to a given output token, and test whether this metric differs systematically between matched surface recall and insight prompts.

Our contributions are:
\begin{enumerate}
    \item A documented negative result: surface text statistics cannot reliably distinguish matched surface/insight pairs after rigorous vocabulary control, across five experimental iterations.
    \item A computational-graph alternative: mean SAE feature layer depth as a proxy for processing depth, tested via circuit tracing on Gemma-2-2B.
    \item A pre-registered replication: the directional hypothesis (surface $>$ insight in mean layer depth) is confirmed on 30 held-out pairs ($p{=}0.000051$, LOO robust).
    \item A robustness analysis demonstrating the depth effect is orthogonal to prediction confidence, ruling out task difficulty as a confound.
\end{enumerate}

%% ============================================================
\section{Related Work}
%% ============================================================

\paragraph{Probing classifiers.}
Linear probing has become a standard tool for investigating what information is encoded in neural network representations \citep{belinkov2022probing, hewitt2019structural}. While effective for detecting syntactic and semantic features, probing is known to be vulnerable to surface confounds \citep{hewitt2019designing}. Our negative result in Section~\ref{sec:negative} provides a concrete case study of this limitation for the surface/insight distinction.

\paragraph{Sparse autoencoders and feature discovery.}
Sparse autoencoders trained on language model activations have been shown to recover interpretable features that correspond to human-understandable concepts \citep{cunningham2023sparse, bricken2023monosemanticity}. Recent work has scaled these methods to production models, identifying features corresponding to factual recall, planning, and multi-step reasoning \citep{anthropic2025biology}.

\paragraph{Circuit tracing and computational graphs.}
Circuit tracing attributes model outputs to specific features and their interactions, enabling analysis of computational pathways rather than static representations \citep{anthropic2025biology}. This approach revealed that multi-hop reasoning in language models proceeds through intermediate computational steps that are invisible in the output text. Our work builds on this methodology, using the \textit{distribution} of feature activations across layers as a summary statistic for processing depth.

\paragraph{Reasoning depth in language models.}
Prior work has explored whether language models exhibit qualitatively different processing for different types of queries \citep{wei2022chain, dziri2023faith}, but typically at the behavioral level (output quality) rather than the mechanistic level (internal computation). Our approach bridges this gap by directly measuring computational depth via SAE attribution.

%% ============================================================
\section{Methods}
%% ============================================================

\subsection{Prompt Pair Construction}

We construct matched prompt pairs where each pair consists of a \textit{surface} prompt (designed to elicit factual recall) and an \textit{insight} prompt (designed to elicit explanatory reasoning) on the same topic. Pairs are matched on the following criteria:

\begin{itemize}
    \item \textbf{Vocabulary:} both prompts in a pair use the same word set. No new content words are introduced in one variant that are absent in the other.
    \item \textbf{Token count:} both prompts have equal or near-equal token counts ($\pm 2$ tokens).
    \item \textbf{Syntactic structure:} sentence structure is kept parallel where possible, varying only the semantic content (factual statement vs.\ explanatory mechanism).
    \item \textbf{Single-token targets:} both prompts target a single next token for attribution, verified prior to each experimental run.
\end{itemize}

Topics were drawn from diverse domains (biology, physics, economics, psychology, history, technology) and generated using Claude 3.5 Sonnet with vocabulary-enforcement constraints. Pairs failing the vocabulary check were regenerated up to three times; pairs that could not be made clean were excluded.

\subsection{Surface Statistics Protocol}
\label{sec:surface_stats}

Prior to the circuit tracing experiments, we conducted five iterations of surface statistics analysis to determine whether the surface/insight distinction could be recovered from text features alone:

\begin{enumerate}
    \item \textbf{Iteration 1:} Bag-of-words TF-IDF classification. Result: $>$84\% accuracy, indicating severe vocabulary confound.
    \item \textbf{Iteration 2:} Vocabulary-controlled pairs with TF-IDF baseline. Result: accuracy remained above chance due to word \textit{frequency} shifts even with identical vocabularies.
    \item \textbf{Iteration 3:} Readability metrics (Flesch-Kincaid, Coleman-Liau, Gunning fog). Result: not reliably discriminative after vocabulary control.
    \item \textbf{Iteration 4:} N-gram overlap analysis. Result: at chance after vocabulary matching.
    \item \textbf{Iteration 5:} Lexical diversity metrics. Result: at chance after vocabulary matching.
\end{enumerate}

The consistent pattern: surface statistics either discriminated the pairs (indicating confounds) or fell to chance (after confound elimination). This motivated the computational approach.

\subsection{Circuit Tracing and Feature Attribution}
\label{sec:circuit_tracing}

For each prompt, we run circuit tracing using the \texttt{circuit-tracer} library on Gemma-2-2B with pre-trained SAE dictionaries. The procedure is:

\begin{enumerate}
    \item Load the prompt and identify the target token (the single next token to be predicted).
    \item Run forward attribution, collecting all SAE features that are attributed to the target token above a minimum activation threshold.
    \item For each attributed feature, record its layer index (0--25 for Gemma-2-2B's 26 layers).
    \item Compute the \textbf{mean layer index} across all attributed features, weighted by attribution magnitude.
\end{enumerate}

The mean layer index serves as our proxy for processing depth. The hypothesis is that surface recall prompts concentrate attribution at higher (output-proximal) layers, reflecting direct retrieval, while insight prompts engage earlier layers, reflecting distributed, multi-step computation.

\subsection{Pre-registration and Replication Design}

The experimental program proceeded in three phases:

\begin{enumerate}
    \item \textbf{Exploratory phase} ($n{=}26$ pairs): hypothesis generation. The direction of the effect (surface $>$ insight in mean layer index) was observed but not pre-registered.
    \item \textbf{Initial replication} ($n{=}15$ held-out pairs, 11 with valid attributions): pre-registered directional hypothesis. This replication was directionally consistent but statistically underpowered.
    \item \textbf{Power replication} ($n{=}30$ new pairs, zero topic overlap with prior phases): pre-registered directional hypothesis with adequate statistical power. All targets were verified as single-token prior to the run.
\end{enumerate}

The one-tailed test is justified by the pre-registered directional hypothesis: we predicted surface $>$ insight in mean layer depth before observing the power replication data.

%% ============================================================
\section{Results}
%% ============================================================

\subsection{Negative Result: Surface Statistics}
\label{sec:negative}

Table~\ref{tab:surface_stats} summarizes the surface statistics results across five iterations. After rigorous vocabulary control, no surface text metric reliably discriminated surface from insight prompts.

\begin{table}[h]
\centering
\caption{Surface statistics discrimination accuracy after vocabulary control. Chance level is 50\%.}
\label{tab:surface_stats}
\begin{tabular}{lcc}
\toprule
\textbf{Metric} & \textbf{Before Control} & \textbf{After Control} \\
\midrule
TF-IDF (bag-of-words) & 84\% & $\sim$50\% \\
Readability (Flesch-Kincaid) & Variable & At chance \\
N-gram overlap & Variable & At chance \\
Lexical diversity & Variable & At chance \\
\bottomrule
\end{tabular}
\end{table}

This establishes that the surface/insight distinction, once vocabulary and distributional confounds are removed, is not detectable from input text statistics. The information, if it exists, must be in the model's internal computation.

\subsection{Exploratory Phase}

On 26 vocabulary-controlled pairs, mean layer index was higher for surface prompts than insight prompts with a mean delta of $+0.26$. The effect was observed in 80\% of pairs ($\approx$21/26). This phase was exploratory: the direction of the effect was not predicted in advance. Based on the observed direction, we formulated the hypothesis: surface recall prompts will produce higher mean feature layer indices than matched insight prompts.

\subsection{Initial Replication}
\label{sec:initial_rep}

Fifteen new pairs were constructed on held-out topics. Four pairs were excluded due to attribution extraction failure (at least one prompt in the pair returned zero active features). This exclusion criterion is technical, not directional---the pairs could not be scored, not that they contradicted the hypothesis.

On the 11 valid pairs: $\Delta{=}+0.118$, $t{=}2.04$, $p{=}0.035$ (one-tailed). Eight of 11 pairs (72.7\%) supported the predicted direction. The attenuated effect size relative to the exploratory phase ($0.118$ vs.\ $0.26$) is consistent with typical exploratory-to-confirmatory shrinkage.

However, leave-one-out analysis revealed statistical fragility: 7 of 11 pairs, when individually removed, caused $p$ to exceed 0.05. This is expected at $n{=}11$---small samples are inherently sensitive to individual observations---but it precluded a strong confirmatory claim and motivated the power replication.

\subsection{Power Replication}
\label{sec:power_rep}

Thirty new pairs were constructed on fresh topics with zero overlap to the 41 pairs from prior phases. All 60 target tokens were verified as single-token before the run began.

\begin{table}[h]
\centering
\caption{Power replication results ($n{=}30$).}
\label{tab:power_rep}
\begin{tabular}{lc}
\toprule
\textbf{Statistic} & \textbf{Value} \\
\midrule
$n$ (valid pairs) & 30 \\
Exclusions & 0 \\
Surface mean layer & 8.359 \\
Insight mean layer & 8.102 \\
$\Delta$ (surface $-$ insight) & $+0.257$ \\
Directional support & 83.3\% (25/30) \\
$t$-statistic & 4.500 \\
$p$-value (one-tailed) & 0.000051 \\
LOO robust & Yes (survives removal of any single pair) \\
\bottomrule
\end{tabular}
\end{table}

The effect size ($+0.257$) is nearly identical to the exploratory estimate ($+0.26$), indicating that the exploratory phase provided an accurate estimate and the initial replication's attenuated effect was due to low power, not a weaker true effect. Full leave-one-out robustness---the result survives removal of any individual pair---addresses the fragility concern from the initial replication.

\subsection{Robustness: Prediction Confidence Confound Analysis}
\label{sec:confound}

A potential confound is that the depth effect could be driven by prediction difficulty rather than the semantic distinction between surface recall and explanatory reasoning. If surface prompts have more predictable (higher-confidence) next tokens, the model might allocate computation differently for difficulty reasons unrelated to the surface/insight distinction.

To test this, we computed the log-probability of the target token under the model for each prompt, yielding a per-pair confidence delta (surface $-$ insight). We then tested whether this confidence delta correlated with the depth delta.

\begin{table}[h]
\centering
\caption{Prediction confidence confound analysis.}
\label{tab:confound}
\begin{tabular}{lcc}
\toprule
\textbf{Test} & \textbf{Statistic} & \textbf{$p$-value} \\
\midrule
Pearson $r$(confidence, depth) & $0.001$ & $0.996$ \\
Spearman $\rho$(confidence, depth) & $-0.008$ & $0.966$ \\
\midrule
\multicolumn{3}{l}{\textit{Descriptive statistics}} \\
\midrule
Surface more confident & 60\% of pairs & --- \\
Mean LP delta (surface $-$ insight) & $+1.97$ & $0.051$ \\
\bottomrule
\end{tabular}
\end{table}

Surface prompts are marginally more confident on average (60\% of pairs show higher log-probability for the surface target, mean LP delta $= +1.97$, $p = 0.051$). This is expected: factual prompts should have more predictable next tokens. Critically, however, this confidence difference has \textit{zero} relationship with the depth effect. The Pearson correlation between confidence delta and depth delta is $r = 0.001$ ($p = 0.996$); the Spearman rank correlation is $\rho = -0.008$ ($p = 0.966$).

The pairs where surface prompts are substantially more confident than their insight counterparts do not show larger depth differences than pairs where confidence is similar. The depth signal is purely semantic---driven by the category of the prompt (surface vs.\ insight), not by how difficult the next-token prediction is.

This orthogonality between prediction confidence and processing depth rules out task difficulty as a confound and strengthens the interpretation that the depth effect reflects a genuine computational distinction between factual retrieval and explanatory reasoning.

\subsection{Preliminary Finding: Cross-System Universality}
\label{sec:universality}

As an exploratory extension, we compared the layer activation topology observed in Gemma-2-2B to published human fMRI activation patterns for analogous cognitive tasks. Neuroscience literature reports that recall tasks concentrate activation in hippocampal and output-proximal regions, while reasoning tasks engage distributed prefrontal processing \citep{miller2001integrative}. We normalized both transformer and fMRI activation patterns to a common 3-vector representation $[\text{early\_ratio}, \text{mid\_ratio}, \text{late\_ratio}]$ and computed a universality coefficient defined as $1 - \bar{\sigma}^2$ across sources, where $\bar{\sigma}^2$ is the mean variance of corresponding signature components.

Across three sources (Gemma-2-2B surface/insight signatures and human fMRI reference data from reasoning and recall task conditions), the universality coefficient was $0.830$ (surface signature variance: $0.031$; insight signature variance: $0.003$). While the fMRI reference data is derived from published group-level summaries rather than individual-subject measurements, the convergence is notable: both systems show early-distributed activation for reasoning and late-concentrated activation for recall. This preliminary result does not establish causality or shared mechanism, but it suggests the depth topology identified in Sections~\ref{sec:power_rep}--\ref{sec:confound} may reflect a general property of information processing systems that differentiate retrieval from reasoning, rather than a transformer-specific artifact.

%% ============================================================
\section{Discussion}
%% ============================================================

\paragraph{What the results show.}
Feature activation depth in SAE circuit representations distinguishes surface recall from explanatory reasoning in Gemma-2-2B. Surface recall concentrates feature attribution at output-proximal layers; explanatory reasoning engages a more distributed pattern with greater earlier-layer involvement. This effect is robust to pre-registration, leave-one-out sensitivity analysis, and is orthogonal to prediction confidence.

\paragraph{What the results do not show.}
We do not claim that mean layer depth is a general-purpose reasoning quality metric. The finding is specific to: (a) Gemma-2-2B, (b) vocabulary-controlled matched pairs, (c) the surface/insight binary distinction. Generalization to larger models, different architectures, uncontrolled naturalistic text, or finer-grained reasoning taxonomies remains untested.

\paragraph{The methodology as contribution.}
Beyond the specific finding, we argue the methodology has independent value. The combination of matched-pair construction with rigorous vocabulary control, SAE feature attribution via circuit tracing, and mean layer depth as a summary statistic provides a reusable framework for investigating computational distinctions in language models. The five-iteration negative result on surface statistics (Section~\ref{sec:negative}) serves as a cautionary baseline: any future study claiming to detect reasoning depth from text features must first demonstrate that it is not detecting vocabulary.

\paragraph{Connection to prior work.}
Our finding is consistent with \citet{anthropic2025biology}'s observation that multi-step reasoning in language models proceeds through intermediate computational steps at earlier layers before converging to output-layer features. The surface/insight distinction may correspond to the difference between single-hop retrieval (short computational path, output-proximal) and multi-hop reasoning (longer path, distributed activation).

\paragraph{Limitations.}
The study uses a single model (Gemma-2-2B, 26 layers) and a total of 56 matched pairs across all phases. The SAE dictionaries and circuit tracing infrastructure constrain the analysis to models with available sparse autoencoders. The vocabulary-matching procedure, while rigorous, may introduce subtle biases in prompt naturality. Finally, ``surface recall'' and ``explanatory reasoning'' are operationalized through prompt design rather than ground-truth cognitive taxonomy, and the boundary between them is not always sharp.

\paragraph{Open Questions and Future Directions.}
Four tractable extensions follow directly from this methodology. \emph{(i) Scale:} Does the effect generalize beyond Gemma-2-2B to 7B, 70B, and frontier-scale models? The matched-pair protocol applies directly to any model with available SAEs. \emph{(ii) Naturalistic text:} The current protocol uses controlled sentence completions; whether the effect holds on uncontrolled naturalistic prompts is an open empirical question. \emph{(iii) Architecture:} Is the effect Gemma-specific, or does it replicate across Llama, Mistral, and other families with available circuit-tracing infrastructure?

\emph{(iv) Depth recovery after interruption.} This is arguably the most discriminating extension. The question ``does reasoning depth increase over time as an agent accumulates context?'' presupposes continuity. Deployed agents face API timeouts, context resets, and session interruptions as operational realities. A more precise formulation: measure an agent's depth profile before a context interruption and again after re-initialization on the same task. An agent backed by genuine distributed LLM reasoning re-establishes a characteristic depth distribution as context rebuilds. A scripted or shallow-retrieval agent never had one. The depth delta at interruption boundary is therefore a clean discriminator between depth that is \emph{architectural} (persistent across resets) versus \emph{contextual} (dependent on accumulated context). We refer to this measurement pipeline as \textbf{Fathom}, and identify the depth-recovery protocol as its primary live-evaluation use case in multi-agent systems.

\paragraph{Scope: Reasoning Type vs Answer Correctness.}
Three experiments tested whether depth score predicts answer correctness on TruthfulQA ($n{=}140$ total pairs). Results were null across all designs ($p{=}0.62$ in the most rigorous test). We conclude that Fathom measures reasoning \emph{type}---whether computation is distributed across early layers or concentrated at output layers---but not whether the answer produced is factually correct. These are distinct constructs. Surface recall of correct facts produces shallow signatures; incorrect answers generate similar shallow signatures. Depth is orthogonal to accuracy. This sharpens rather than weakens the core claim.

\paragraph{Cognitive Fingerprinting and Causal Intervention.}
The Fathom pipeline has been extended beyond mean layer depth to capture the full shape of a model's computation. A \emph{cognitive fingerprint} represents each prompt's processing signature as a fixed-length vector: $[\text{weighted\_depth}, \text{layer\_entropy}, \text{early\_ratio}, \text{mid\_ratio}, \text{late\_ratio}, \text{feature\_engagement}, \text{compute\_efficiency}]$, enabling clustering and downstream classification of reasoning patterns. A \emph{causal intervention} module performs layer-wise ablation---zeroing SAE features at each layer and measuring the resulting output probability delta---to distinguish layers that are merely active from those that are \emph{causally necessary} for the model's output. This yields a causal depth score in which layers are weighted by functional importance rather than feature count. Finally, a \emph{phase transition detector} applies change-point analysis to token-level depth trajectories, identifying the moment during generation when a model commits to deep processing versus surface recall. These extensions transform the single-number depth metric into a multi-dimensional cognitive profile, and their validation on held-out data is ongoing.

%% ============================================================
\section{Conclusion}
%% ============================================================

We show that sparse autoencoder feature activation depth distinguishes surface recall from explanatory reasoning in Gemma-2-2B, where surface text statistics cannot. The effect is pre-registered, replicated at $p < 0.001$ with full leave-one-out robustness, and orthogonal to prediction confidence. These results demonstrate that computational graph analysis can detect distinctions in model processing that are invisible to representational probing, and provide a reusable methodology for investigating reasoning depth in language models.

%% ============================================================

\bibliographystyle{plainnat}

\begin{thebibliography}{10}

\bibitem[Anthropic(2025)]{anthropic2025biology}
Anthropic.
\newblock On the biology of a large language model.
\newblock Technical report, Anthropic, 2025.

\bibitem[Belinkov(2022)]{belinkov2022probing}
Yonatan Belinkov.
\newblock Probing classifiers: Promises, shortcomings, and advances.
\newblock \textit{Computational Linguistics}, 48(1):207--219, 2022.

\bibitem[Bricken et~al.(2023)]{bricken2023monosemanticity}
Trenton Bricken, Adly Templeton, Joshua Batson, Brian Chen, Adam Jermyn, Tom Conerly, Nick Turner, Cem Anil, Carson Denison, Amanda Askell, Robert Lasenby, Yifan Wu, Shauna Kravec, Nicholas Schiefer, Tim Maxwell, Nicholas Joseph, Zac Hatfield-Dodds, Alex Tamkin, Karina Nguyen, Brayden McLean, Josiah~E Burke, Tristan Hume, Shan Carter, Tom Henighan, and Christopher Olah.
\newblock Towards monosemanticity: Decomposing language models with dictionary learning.
\newblock \textit{Transformer Circuits Thread}, 2023.

\bibitem[Cunningham et~al.(2023)]{cunningham2023sparse}
Hoagy Cunningham, Aidan Ewart, Logan Riggs, Robert Huben, and Lee Sharkey.
\newblock Sparse autoencoders find highly interpretable features in language models.
\newblock \textit{arXiv preprint arXiv:2309.08600}, 2023.

\bibitem[Dziri et~al.(2023)]{dziri2023faith}
Nouha Dziri, Ximing Lu, Melanie Sclar, Xiang~Lorraine Li, Liwei Jiang, Bill~Yuchen Lin, Sean Welleck, Peter West, Chandra Bhagavatula, Ronan Le~Bras, Jena~D Hwang, Soumya Sanyal, Sean Yang, Xiang Ren, Allyson Ettinger, Za\"id Harchaoui, and Yejin Choi.
\newblock Faith and fate: Limits of transformers on compositionality.
\newblock In \textit{NeurIPS}, 2023.

\bibitem[Hewitt and Liang(2019)]{hewitt2019designing}
John Hewitt and Percy Liang.
\newblock Designing and interpreting probes with control tasks.
\newblock In \textit{EMNLP}, 2019.

\bibitem[Hewitt and Manning(2019)]{hewitt2019structural}
John Hewitt and Christopher~D Manning.
\newblock A structural probe for finding syntax in word representations.
\newblock In \textit{NAACL}, 2019.

\bibitem[Miller and Cohen(2001)]{miller2001integrative}
Earl~K Miller and Jonathan~D Cohen.
\newblock An integrative theory of prefrontal cortex function.
\newblock \textit{Annual Review of Neuroscience}, 24:167--202, 2001.

\bibitem[Wei et~al.(2022)]{wei2022chain}
Jason Wei, Xuezhi Wang, Dale Schuurmans, Maarten Bosma, Brian Ichter, Fei Xia, Ed~Chi, Quoc~V Le, and Denny Zhou.
\newblock Chain-of-thought prompting elicits reasoning in large language models.
\newblock In \textit{NeurIPS}, 2022.

\end{thebibliography}

\end{document}
